% ============================================================
%  AttentiON — Memoria TFG
%  Autor: Unai Fernández
%  Grado: Diseño Digital y Tecnologías Multimedia
%  Universidad: Centro de la Imagen y Tecnología Multimedia
%  Curso: 2024-25
% ============================================================

\documentclass[12pt, a4paper]{report}

% ─── Codificación y lengua ───────────────────────────────────
\usepackage[utf8]{inputenc}
\usepackage[T1]{fontenc}
\usepackage[spanish]{babel}

% ─── Tipografía ──────────────────────────────────────────────
\usepackage{lmodern}

% ─── Geometría ───────────────────────────────────────────────
\usepackage[
  top=2.5cm, bottom=2.5cm,
  left=3cm,  right=2.5cm
]{geometry}

% ─── Interlineado ────────────────────────────────────────────
\usepackage{setspace}
\onehalfspacing

% ─── Colores y gráficos ──────────────────────────────────────
\usepackage{xcolor}
\usepackage{graphicx}
\usepackage{float}

\definecolor{accentblue}{RGB}{45, 120, 210}
\definecolor{darkgray}{RGB}{60,60,60}
\definecolor{lightgray}{RGB}{240,240,240}
\definecolor{codebg}{RGB}{248,248,252}

% ─── Tablas ──────────────────────────────────────────────────
\usepackage{booktabs}
\usepackage{tabularx}
\usepackage{multirow}
\usepackage{longtable}
\usepackage{array}

% ─── Listas ──────────────────────────────────────────────────
\usepackage{enumitem}

% ─── Código fuente ───────────────────────────────────────────
\usepackage{listings}
\usepackage{inconsolata}

\lstdefinelanguage{CSharp}{
  morekeywords={abstract, as, base, bool, break, byte, case, catch, char,
    checked, class, const, continue, decimal, default, delegate, do, double,
    else, enum, event, explicit, extern, false, finally, fixed, float, for,
    foreach, goto, if, implicit, in, int, interface, internal, is, lock, long,
    namespace, new, null, object, operator, out, override, params, private,
    protected, public, readonly, ref, return, sbyte, sealed, short, sizeof,
    stackalloc, static, string, struct, switch, this, throw, true, try, typeof,
    uint, ulong, unchecked, unsafe, ushort, using, virtual, void, volatile,
    while, get, set, var, yield, async, await, partial, dynamic},
  sensitive=true,
  morecomment=[l]{//},
  morecomment=[s]{/*}{*/},
  morestring=[b]",
}

\lstset{
  language=CSharp,
  backgroundcolor=\color{codebg},
  basicstyle=\footnotesize\ttfamily,
  keywordstyle=\color{accentblue}\bfseries,
  stringstyle=\color{darkgray},
  commentstyle=\color{gray}\itshape,
  numberstyle=\tiny\color{gray},
  numbers=left,
  numbersep=8pt,
  frame=single,
  framerule=0.4pt,
  rulecolor=\color{gray!40},
  breaklines=true,
  breakatwhitespace=false,
  tabsize=4,
  showstringspaces=false,
  xleftmargin=1.5em,
  captionpos=b,
}

% ─── Hiperenlaces ────────────────────────────────────────────
\usepackage[
  colorlinks=true,
  linkcolor=accentblue,
  urlcolor=accentblue,
  citecolor=accentblue,
  bookmarks=true,
  bookmarksopen=true,
  pdftitle={AttentiON - Memoria TFG},
  pdfauthor={Unai Fernández}
]{hyperref}

% ─── Encabezados y pies ──────────────────────────────────────
\usepackage{fancyhdr}
\pagestyle{fancy}
\fancyhf{}
\fancyhead[L]{\small\textcolor{darkgray}{AttentiON}}
\fancyhead[R]{\small\textcolor{darkgray}{Unai Fernández}}
\fancyfoot[C]{\thepage}
\renewcommand{\headrulewidth}{0.4pt}

% ─── Formato de capítulos ────────────────────────────────────
\usepackage{titlesec}
\titleformat{\chapter}[hang]
  {\Large\bfseries\color{accentblue}}
  {\thechapter.}{0.8em}{}
  [\vspace{0.3em}\hrule]
\titlespacing*{\chapter}{0pt}{10pt}{20pt}

\titleformat{\section}
  {\large\bfseries\color{darkgray}}
  {\thesection}{0.7em}{}

\titleformat{\subsection}
  {\normalsize\bfseries\color{darkgray}}
  {\thesubsection}{0.6em}{}

% ─── Bibliografía ────────────────────────────────────────────
\usepackage[authoryear, round]{natbib}

% ─── Miscelánea ──────────────────────────────────────────────
\setlength{\parindent}{0pt}
\setlength{\parskip}{6pt}

% ============================================================
%  DOCUMENTO
% ============================================================
\begin{document}

\pagenumbering{roman}

% ─── PORTADA ─────────────────────────────────────────────────
\begin{titlepage}
  \centering
  \vspace*{2cm}

  {\Huge\bfseries\color{accentblue} AttentiON}\\[0.5cm]
  {\large Plataforma de minijuegos educativos para el desarrollo cognitivo}\\[2cm]

  {\large\textbf{Unai Fernández}}\\[1.5cm]

  \begin{tabular}{rl}
    \textbf{Director:} & Unai Fernández \\[4pt]
    \textbf{Grado:}    & Diseño Digital y Tecnologías Multimedia \\[4pt]
    \textbf{Curso:}    & 2024--25 \\[4pt]
    \textbf{Universidad:} & Centro de la Imagen y Tecnología Multimedia \\
  \end{tabular}

  \vfill

  {\small\textcolor{gray}{Trabajo Final de Grado — Curso 2024-25}}
\end{titlepage}

% ─── RESUMEN ─────────────────────────────────────────────────
\chapter*{Resumen}
\addcontentsline{toc}{chapter}{Resumen}

La presente memoria describe el desarrollo del Trabajo Final de Grado
\textit{AttentiON}, una plataforma de videojuego educativo implementada con el
motor Unity (C\#) orientada al entrenamiento cognitivo de niños y niñas en
edad escolar (7--12 años).

El proyecto parte de la necesidad de disponer de recursos digitales interactivos
que combinen valor pedagógico y atractivo lúdico, superando las limitaciones
detectadas en las aplicaciones existentes del sector. Con ese fin, se ha diseñado
y desarrollado un sistema modular de minijuegos clasificados en cinco categorías
de habilidades ejecutivas: atención, memoria, control de impulsos, gestión
emocional y planificación.

Desde el punto de vista técnico, la plataforma se articula sobre una arquitectura
basada en el patrón Singleton para los gestores centrales
(\texttt{GameManager}, \texttt{UIManager}) y una clase base abstracta
(\texttt{MinigameBase}) que unifica el ciclo de vida de los minijuegos. Se han
implementado un sistema de dificultad adaptativa con tres niveles (Fácil,
Medio, Difícil), un sistema de puntuación acumulativa, un sistema de carga de
escenas centralizado y un panel de introducción generado proceduralmente para
cada minijuego.

La metodología de trabajo ha sido iterativa e incremental, lo que ha permitido
desarrollar, probar y pulir cada minijuego de manera independiente antes de su
integración final. El prototipo resultante incluye más de veinte minijuegos
funcionales distribuidos entre las cinco categorías, con interfaces adaptadas
al público infantil y una estética visual coherente.

\bigskip
\textbf{Palabras clave:} videojuego educativo, serious games, atención, memoria operativa,
control de impulsos, habilidades ejecutivas, minijuego, gamificación, Unity, C\#,
experiencia de usuario (UX).

% ─── ENLACES ─────────────────────────────────────────────────
\chapter*{Enlaces}
\addcontentsline{toc}{chapter}{Enlaces}

A continuación se recogen los recursos digitales creados específicamente en el
marco de este proyecto:

\begin{itemize}
  \item \textbf{Repositorio del proyecto (GitHub):}\\
    \url{https://github.com/} \\
    Código fuente del proyecto Unity, historial de versiones y documentación
    técnica.

  \item \textbf{Diagrama de Gantt:}\\
    \url{https://docs.google.com/spreadsheets/d/19SBDgHSgRw6bz_plAlHYkfWy23hx8tOTT-GIoEdbngE/edit?usp=sharing}\\
    Planificación temporal del proyecto con fases, tareas y fechas de entrega.
\end{itemize}

% ─── ÍNDICE DE TABLAS ─────────────────────────────────────────
\listoftables
\addcontentsline{toc}{chapter}{Índice de tablas}

% ─── ÍNDICE DE FIGURAS ────────────────────────────────────────
\listoffigures
\addcontentsline{toc}{chapter}{Índice de figuras}

% ─── GLOSARIO ─────────────────────────────────────────────────
\chapter*{Glosario}
\addcontentsline{toc}{chapter}{Glosario}

% (El glosario se completará en la versión final del documento.)

% ─── ÍNDICE GENERAL ───────────────────────────────────────────
\tableofcontents

\clearpage
\pagenumbering{arabic}

% ============================================================
%  CAPÍTULO 1 — INTRODUCCIÓN
% ============================================================
\chapter{Introducción}

La presente memoria desarrolla el Trabajo Final de Grado \textit{AttentiON},
un proyecto orientado a explorar el potencial educativo de los videojuegos
en el desarrollo de habilidades cognitivas en edades tempranas. En los
siguientes subapartados se contextualizan la motivación académica, el
problema detectado, los objetivos generales y específicos, y las posibilidades
del prototipo resultante.

\section{Motivación}

El auge de la tecnología digital en el ámbito educativo ha abierto nuevas
posibilidades para el diseño de herramientas de aprendizaje. Los videojuegos,
en particular, han demostrado ser entornos especialmente eficaces para estimular
habilidades cognitivas de forma lúdica y atractiva \citep{Gee2003,Connolly2012}.

El proyecto \textit{AttentiON} nace de la necesidad de proponer modalidades de
aprendizaje renovadas, fundamentadas en la interacción, el entretenimiento y la
tecnología. La motivación central es la creación de un entorno digital que permita
a niños y jóvenes trabajar aspectos esenciales como la atención sostenida, la
memoria operativa y el control de impulsos, todo ello desde una perspectiva lúdica
y sumamente motivadora.

Además del componente tecnológico, el proyecto está impulsado por la convicción de
que el aprendizaje puede trascender los métodos pedagógicos más convencionales,
integrando el juego como vehículo para el desarrollo cognitivo y emocional
\citep{Granic2014}. Esta perspectiva no es exclusivamente personal: existe una
base científica sólida que avala el uso de los videojuegos como recurso educativo
eficaz \citep{Bediou2018,Powers2013}.

\section{Formulación del problema}

En el mercado actual existen diversas plataformas de entrenamiento cognitivo
digital, como Lumosity, CogniFit o BrainHQ. Sin embargo, estas aplicaciones
presentan limitaciones significativas cuando se trata de adaptarlas al público
infantil: las interfaces tienden a resultar poco atractivas para los más jóvenes,
y las mecánicas se vuelven repetitivas con rapidez, lo que provoca una disminución
del interés y del compromiso del usuario \citep{Wouters2013,Boyle2016}.

El reto central del proyecto es, por tanto, doble: por un lado, diseñar ejercicios
que estimulen funciones cognitivas relevantes; por otro, envolverlos en una
propuesta visual moderna y una jugabilidad lo suficientemente entretenida para
que el jugador no perciba la actividad como una tarea escolar más. No se trata
únicamente de mejorar soluciones existentes, sino de plantear una propuesta
diferente que integre criterios pedagógicos con principios de diseño de
videojuegos actuales.

\section{Objetivos generales del TFG}

El objetivo principal de este Trabajo Final de Grado es desarrollar un prototipo
funcional de videojuego educativo orientado al entrenamiento de las funciones
ejecutivas mediante dinámicas interactivas y entretenidas.

El juego está concebido para ofrecer una experiencia atractiva, sencilla de
comprender y accesible para niños en edad escolar, de modo que puedan
beneficiarse de sus mecánicas cognitivas sin necesidad de formación previa
ni supervisión constante.

\section{Objetivos específicos del TFG}

\begin{itemize}
  \item Diseñar un sistema de menús y una interfaz intuitiva adaptada al
    público infantil.
  \item Implementar minijuegos funcionales centrados en habilidades cognitivas
    específicas: memoria operativa, atención sostenida, control de impulsos,
    planificación y gestión emocional.
  \item Incorporar una estética visual vibrante y coherente con la temática
    educativa.
  \item Estructurar el contenido en niveles de dificultad progresivos
    (Fácil, Medio, Difícil) adaptados a diferentes franjas de edad.
  \item Evaluar la jugabilidad mediante pruebas internas del prototipo.
  \item Documentar el proceso completo de diseño, desarrollo y justificación
    pedagógica del videojuego.
\end{itemize}

\section{Alcance del proyecto}

El alcance del proyecto comprende el desarrollo de un prototipo jugable
funcional, concebido como una primera versión de una plataforma interactiva
de entrenamiento cognitivo basada en el juego.

El prototipo incluye los siguientes componentes principales:

\begin{itemize}
  \item Un menú inicial operativo desde el cual el jugador puede navegar
    por las distintas secciones del sistema de forma intuitiva.
  \item Un selector de dificultad con tres niveles temáticos (Bosque Mágico,
    Castillo Volador y Cueva Misteriosa), que adapta la experiencia a
    diferentes franjas de edad.
  \item Un módulo de selección de minijuegos organizado en cinco categorías
    cognitivas: atención, memoria, control de impulsos, gestión emocional y
    planificación.
  \item Una serie de minijuegos cortos e independientes, cada uno centrado
    en una habilidad ejecutiva concreta. La brevedad de cada sesión es un
    principio de diseño clave: la experiencia debe ser motivadora y no
    generar fatiga cognitiva.
\end{itemize}

El proyecto se dirige principalmente a niños y niñas en edad escolar, con
edades comprendidas entre los 7 y los 12 años, franja en la que el desarrollo
cognitivo es especialmente receptivo a los aprendizajes mediante dinámicas
lúdicas. No obstante, la plataforma puede resultar igualmente útil para
educadores y psicopedagogos que busquen herramientas interactivas para
complementar sus metodologías de enseñanza.

Los beneficiarios directos son los propios jugadores, que encontrarán en la
plataforma una forma divertida y significativa de fortalecer habilidades
cognitivas fundamentales. El impacto se extiende, además, a maestros,
psicopedagogos, instituciones educativas y familias, quienes podrán utilizarla
como apoyo en procesos de aprendizaje, estimulación cognitiva y detección
temprana de dificultades.

Quedan fuera del alcance del presente trabajo la validación científica clínica
de los resultados cognitivos, la publicación comercial de la plataforma y el
desarrollo de un sistema de seguimiento del progreso del jugador a largo plazo.
Estas líneas de trabajo se identifican como posibles extensiones futuras.

% ============================================================
%  CAPÍTULO 2 — ESTADO DEL ARTE
% ============================================================
\chapter{Estado del arte / Marco teórico / Contextualización / Estudio de mercado}

Este capítulo analiza los antecedentes teóricos y prácticos relacionados con
los videojuegos educativos, su impacto en las funciones cognitivas y el estado
actual del mercado de los serious games. Se revisan estudios científicos
relevantes y se comparan aplicaciones existentes destinadas al entrenamiento
cognitivo, con el fin de identificar las oportunidades que justifican la
propuesta de \textit{AttentiON}.

\section{Definición y origen de los serious games}

El concepto de \textit{serious games} hace referencia a videojuegos, simulaciones
interactivas o experiencias digitales que, aunque pueden incorporar elementos
lúdicos para aumentar la motivación y la atención del jugador, tienen como
objetivo principal algo que va más allá del entretenimiento.

La consolidación del término como disciplina académica se produjo en la primera
década del siglo XXI, impulsada por tres factores convergentes: el ascenso de la
industria del videojuego, el desarrollo de herramientas accesibles para la creación
de contenidos interactivos y el aumento de la investigación que evidenciaba la
eficacia de los videojuegos para el aprendizaje y el entrenamiento
\citep{Michael2006}.

En este sentido, \citet{Michael2006} definen los juegos serios como
«juegos que no tienen el entretenimiento, el disfrute o la diversión como
su propósito principal», y señalan que estos «utilizan el medio artístico de
los juegos para transmitir un mensaje, enseñar una lección o brindar una
experiencia» (traducción propia del original en inglés).

En la actualidad, los \textit{serious games} abarcan desde juegos educativos
en entornos escolares hasta plataformas utilizadas para la rehabilitación
cognitiva o motora, pasando por simuladores complejos aplicados en medicina
o aviación \citep{GamesForChange}.

Los \textit{serious games} se caracterizan por combinar:

\begin{itemize}
  \item Objetivos funcionales claros (educar, entrenar, concienciar, evaluar).
  \item Interactividad significativa, que favorece la toma de decisiones y la
    experimentación.
  \item Narrativas o contextos simulados, que permiten aprender en entornos
    seguros.
  \item Mecánicas motivadoras, que incrementan el compromiso del usuario.
  \item Retroalimentación inmediata, esencial para la adquisición progresiva
    de habilidades.
  \item Medición del rendimiento, útil para entornos educativos, sanitarios
    o de formación profesional.
\end{itemize}

\section{Serious games en el ámbito educativo}

En el entorno educativo, los \textit{serious games} han demostrado ser una
herramienta capaz de ofrecer mucho más que motivación. Al integrar dinámicas
interactivas y situaciones que exigen tomar decisiones en tiempo real, estos
juegos favorecen la concentración y contribuyen a mantener la atención durante
periodos más prolongados, algo que resulta más complicado en el aula tradicional
\citep{Bediou2018}.

Asimismo, permiten desarrollar de forma natural habilidades visuoespaciales,
dado que muchos entornos de juego exigen orientarse, anticipar movimientos o
interpretar información visual compleja \citep{GreenBavelier2003}. Los
retos planteados de forma progresiva fomentan la memoria de trabajo, ya que el
jugador necesita retener datos, compararlos y utilizarlos mientras avanza
\citep{Anguera2013}.

La velocidad de procesamiento también se beneficia: la rapidez con la que se
presentan ciertos estímulos obliga a los jugadores a responder con mayor agilidad.
El control inhibitorio se ve igualmente reforzado, puesto que numerosos juegos
premian la calma, la estrategia y la capacidad de frenar impulsos para reflexionar
antes de actuar \citep{Powers2013}.

La comunidad internacional de \textit{serious games} se articula en torno a
eventos de referencia como el \textit{Games for Change Festival}, la
\textit{Joint Conference on Serious Games} o la \textit{Serious Play Conference},
que reúnen a académicos, investigadores y profesionales de la industria con el
objetivo de avanzar en el campo \citep{GamesForChange,JCSG,SeriousPlay}.

\citet{GamesForChange} describe su propósito en los siguientes términos:
«Los videojuegos son un campo que crece y evoluciona rápidamente, impulsado
por tecnologías emergentes que están transformando la forma en que colaboramos
y nos comunicamos entre nosotros» (traducción propia del original en inglés).

\section{Relación entre videojuegos y habilidades cognitivas}

Diversos estudios han mostrado que ciertos videojuegos pueden influir
positivamente en funciones cognitivas de alto nivel, incluyendo las funciones
ejecutivas \citep{Powers2013,Bediou2018}.

Respecto a la \textbf{memoria de trabajo}, algunas investigaciones sugieren
que los videojuegos diseñados específicamente para ello pueden mejorar esta
capacidad, especialmente cuando la mecánica exige mantener y actualizar
información en paralelo con otras tareas \citep{Anguera2013}.

En cuanto a la \textbf{inhibición}, un metaanálisis constata que los
\textit{serious games} pueden entrenar mejor esta función ejecutiva que los
videojuegos convencionales, al incorporar reglas que obligan a los jugadores
a controlar sus respuestas automáticas para alcanzar los objetivos del juego
\citep{Powers2013}.

La \textbf{atención sostenida y selectiva} es quizás el área donde la evidencia
resulta más sólida: estudios con distintos géneros de videojuegos, incluyendo
MMORPG y \textit{serious games}, han demostrado mejoras significativas en estas
funciones \citep{GreenBavelier2003,Bediou2018}.

Finalmente, la \textbf{gestión emocional} también puede beneficiarse del uso de
videojuegos, especialmente en contextos educativos o terapéuticos. Estudios
recientes señalan que los \textit{serious games} diseñados para trabajar la
regulación emocional tienen un impacto positivo en niños y adolescentes,
ayudándoles a identificar, modular y controlar sus emociones
\citep{Granic2014,UNIR2024}.

\section{Análisis de aplicaciones existentes}

En el campo de los videojuegos educativos orientados al entrenamiento cognitivo
existen varias plataformas de referencia. A continuación se analiza cada una
de ellas atendiendo a sus mecánicas de juego, su base científica y su adecuación
al público infantil.

\subsection{Lumosity}

Lumosity es una de las plataformas de entrenamiento cerebral más extendidas a
nivel mundial. Su enfoque consiste en trasladar experimentos clásicos de
psicología cognitiva a un formato de minijuego interactivo, lo que permite
estimular funciones como la memoria de trabajo, la atención selectiva o la
velocidad de procesamiento de forma controlada y repetitiva.

Lumos Labs llevó a cabo un estudio aleatorizado sobre el entrenamiento cerebral
de Lumosity cuyos resultados fueron publicados en una revista de investigación
revisada por pares. En el estudio, la mitad de los 4.715 participantes que lo
completaron entrenaron cinco días a la semana durante quince minutos al día en
la plataforma, mientras que la otra mitad realizó crucigramas en línea como
grupo de control activo \citep{LumosityWeb}.

Los estudios coinciden en que los usuarios suelen mejorar dentro de los propios
juegos, fenómeno conocido como mejora en la tarea (\textit{task-specific
improvement}). Sin embargo, la evidencia sobre la transferencia de estas mejoras
a situaciones de la vida real es más limitada, especialmente cuando se trata de
habilidades complejas o generales.

\begin{figure}[H]
  \centering
  % \includegraphics[width=0.4\textwidth]{figuras/lumosity.png}
  \fbox{\parbox{6cm}{\centering\vspace{1cm}[Figura: Logotipo de Lumosity]\vspace{1cm}}}
  \caption{Logotipo de Lumosity, plataforma de entrenamiento cognitivo
    mediante minijuegos de memoria, atención y velocidad de procesamiento.}
  \label{fig:lumosity}
\end{figure}

\subsection{CogniFit}

CogniFit es otra de las plataformas más representativas del ámbito del
entrenamiento cognitivo digital. Su funcionamiento se basa en minijuegos
interactivos diseñados para estimular distintas funciones mentales, entre
ellas la memoria de trabajo, la coordinación visomotora, la planificación, la
toma de decisiones y la atención sostenida.

Desde la perspectiva científica, CogniFit fundamenta su propuesta en la
neuropsicología y la plasticidad cerebral: la repetición de tareas puede
fortalecer las redes neuronales asociadas a funciones cognitivas como la
memoria, la atención o la percepción.

En un estudio con diseño aleatorizado controlado de cuatro grupos
(intervención cognitiva, intervención de actividad física, intervención
combinada y grupo de control) se evaluó la eficacia del entrenamiento cognitivo
de la plataforma para mejorar la función cognitiva en adultos mayores sanos
(traducción propia del original en inglés) \citep{CogniFitWeb}.

\begin{figure}[H]
  \centering
  \fbox{\parbox{6cm}{\centering\vspace{1cm}[Figura: Logotipo de CogniFit]\vspace{1cm}}}
  \caption{Logotipo de CogniFit, que representa su enfoque en ejercicios
    personalizados para entrenar habilidades cognitivas específicas.}
  \label{fig:cognifit}
\end{figure}

\subsection{BrainHQ}

BrainHQ es una plataforma de entrenamiento cognitivo desarrollada por Posit
Science en colaboración con neurocientíficos de renombre. Sus ejercicios abarcan
la velocidad de procesamiento, la memoria, la atención y la precisión perceptiva.
Una de sus características diferenciadoras es que la dificultad se ajusta
progresivamente a cada jugador, garantizando así un nivel de reto que mantiene
la atención sin provocar frustración.

La descripción oficial de la plataforma indica: «BrainHQ es tu plataforma en
línea para el entrenamiento cerebral, con ejercicios especializados para la
memoria, la atención, la velocidad de procesamiento, las habilidades sociales,
la toma de decisiones y la orientación» (traducción propia del original en
inglés) \citep{BrainHQWeb}.

\begin{figure}[H]
  \centering
  \fbox{\parbox{6cm}{\centering\vspace{1cm}[Figura: Logotipo de BrainHQ]\vspace{1cm}}}
  \caption{Logotipo de BrainHQ, plataforma centrada en el entrenamiento de
    velocidad de procesamiento, memoria y atención mediante ejercicios adaptativos.}
  \label{fig:brainhq}
\end{figure}

\subsection{NeuroRacer}

NeuroRacer es un videojuego diseñado por el Gazzaley Lab de la Universidad de
California para evaluar y entrenar capacidades cognitivas relacionadas con la
atención y el control cognitivo. Su mecánica principal consiste en una tarea
dual: el jugador debe conducir un vehículo por una carretera llena de curvas
mientras responde simultáneamente a estímulos visuales que aparecen en su
campo de visión.

Los resultados del estudio original fueron publicados en la revista científica
\textit{Nature} y demostraron que la experiencia interactiva puede producir
mejoras cognitivas específicas y medibles, correlacionadas con registros
electroencefalográficos \citep{Anguera2013}.

\begin{figure}[H]
  \centering
  \fbox{\parbox{6cm}{\centering\vspace{1cm}[Figura: Captura de NeuroRacer]\vspace{1cm}}}
  \caption{Captura de la interfaz de NeuroRacer, juego experimental donde el
    usuario conduce un vehículo mientras responde a estímulos visuales,
    entrenando la atención selectiva y el control cognitivo.}
  \label{fig:neuroracer}
\end{figure}

\section{Comparación crítica y oportunidades destacadas}

El análisis comparativo de las plataformas revisadas permite identificar
fortalezas comunes y limitaciones compartidas que señalan oportunidades de
mejora concretas para \textit{AttentiON}.

Lumosity, CogniFit y BrainHQ poseen una base científica sólida y ofrecen
ejercicios que entrenan la atención, la memoria y otras funciones cognitivas.
Sin embargo, su público objetivo parece ser adulto o general, lo que las limita
como herramientas para niños en edad escolar. Ninguna de ellas presenta una
narrativa especialmente diseñada para captar y mantener el interés de un usuario
muy joven. NeuroRacer, por su parte, demuestra que el entrenamiento inmersivo
basado en la multitarea puede mejorar la atención sostenida y la memoria de
trabajo, pero exige condiciones de laboratorio que reducen su aplicabilidad
autónoma en el hogar o el aula.

De este análisis se desprenden las siguientes oportunidades que
\textit{AttentiON} busca aprovechar:

\begin{itemize}
  \item \textbf{Diversas habilidades cognitivas en un solo entorno:}
    memoria operativa, control de impulsos, atención sostenida y selectiva,
    planificación y gestión emocional, integradas en una misma plataforma.
  \item \textbf{Diseño atractivo y accesible:} menús intuitivos, estética
    visual vibrante y dinámicas motivadoras adaptadas a niños de 7 a 12 años.
  \item \textbf{Autonomía en el aprendizaje:} juegos breves que pueden jugarse
    de forma independiente, sin supervisión constante, fomentando la iniciativa
    y la autorregulación.
  \item \textbf{Interactividad y retroalimentación inmediata:} mecanismos
    que mantienen el interés y refuerzan el aprendizaje cognitivo de manera
    lúdica.
\end{itemize}

\section{Conclusiones del estado del arte}

El análisis del estado del arte demuestra que los videojuegos educativos y los
\textit{serious games} pueden mejorar significativamente funciones cognitivas
como la memoria de trabajo, la atención sostenida y selectiva, el control de
impulsos, la planificación y la gestión emocional
\citep{Anguera2013,Connolly2012,Granic2014}.

Las aplicaciones existentes presentan limitaciones importantes: la mayoría está
orientada a adultos, carece de atractivo visual para el público infantil y no
integra múltiples habilidades cognitivas en un mismo entorno. En este contexto,
\textit{AttentiON} surge como una propuesta que busca precisamente dar respuesta
a esas carencias, combinando rigor pedagógico con un diseño lúdico adaptado al
público al que se dirige.

% ============================================================
%  CAPÍTULO 3 — GESTIÓN DEL PROYECTO
% ============================================================
\chapter{Gestión del proyecto}

La gestión del proyecto \textit{AttentiON} ha sido la herramienta que ha
permitido organizar un desarrollo que combina investigación teórica con
programación en Unity. Al tratarse de un trabajo individual, ha sido necesario
gestionar simultáneamente las funciones de diseño, programación y redacción,
lo que ha exigido una planificación rigurosa y una autodisciplina estricta
para evitar cuellos de botella y cumplir con los plazos académicos.

Para visualizar este esfuerzo se elaboró una planificación temporal apoyada
en un diagrama de Gantt, accesible en los enlaces del proyecto, que abarca
desde junio de 2025 hasta junio de 2026. Gracias a esta herramienta ha
sido posible monitorizar el grado de avance de las distintas fases y ajustar
el ritmo de trabajo ante dificultades técnicas imprevistas.

\section{Análisis DAFO}

Con el objetivo de evaluar la viabilidad y el contexto del proyecto, se ha
elaborado un análisis DAFO que identifica los principales factores internos
y externos que influyen en su desarrollo.

\begin{table}[H]
\centering
\caption{Análisis DAFO del proyecto AttentiON}
\label{tab:dafo}
\renewcommand{\arraystretch}{1.3}
\begin{tabularx}{\textwidth}{|>{\bfseries}p{3.5cm}|X|}
\hline
\multicolumn{2}{|c|}{\textbf{FACTORES INTERNOS}} \\
\hline
Fortalezas &
\begin{itemize}[leftmargin=*, nosep, after=\vspace{2pt}]
  \item Propuesta pedagógica claramente definida: el videojuego tiene un
    objetivo explícito de mejora de funciones ejecutivas.
  \item Base teórica y científica sólida, apoyada en psicología cognitiva
    y en estudios sobre \textit{serious games}.
  \item Prototipo funcional desarrollado en Unity, con valor práctico tangible.
  \item Arquitectura modular basada en minijuegos, que facilita la escalabilidad.
  \item Perfil multidisciplinar del autor, que integra diseño visual y
    programación técnica.
\end{itemize} \\
\hline
Debilidades &
\begin{itemize}[leftmargin=*, nosep, after=\vspace{2pt}]
  \item Alcance condicionado por el trabajo individual y los plazos académicos.
  \item Imposibilidad de completar los 25 minijuegos inicialmente planificados.
  \item Limitación de recursos propios del ámbito académico.
  \item Posible disparidad en el nivel de acabado entre minijuegos desarrollados
    en fases distintas.
\end{itemize} \\
\hline
\multicolumn{2}{|c|}{\textbf{FACTORES EXTERNOS}} \\
\hline
Oportunidades &
\begin{itemize}[leftmargin=*, nosep, after=\vspace{2pt}]
  \item Creciente integración de \textit{serious games} en los sistemas escolares.
  \item Demanda social y pedagógica de recursos digitales para el refuerzo
    cognitivo.
  \item Potencial uso como recurso docente complementario en el aula.
  \item Escalabilidad y proyección a largo plazo del prototipo.
  \item Disponibilidad del motor Unity de forma gratuita para proyectos
    académicos.
\end{itemize} \\
\hline
Amenazas &
\begin{itemize}[leftmargin=*, nosep, after=\vspace{2pt}]
  \item Competencia de plataformas ya establecidas con mayores recursos.
  \item Imprevistos en el calendario académico.
  \item Barreras tecnológicas inesperadas en Unity.
  \item Presión por la carga de trabajo acumulada al final del grado.
  \item Dificultad para evaluar con precisión los resultados cognitivos en
    un entorno académico.
\end{itemize} \\
\hline
\end{tabularx}
\end{table}

\section{Riesgos y plan de contingencias}

Durante el desarrollo de \textit{AttentiON} se han identificado los riesgos
más relevantes para el éxito del TFG, clasificándolos según su naturaleza y
estableciendo para cada uno un plan de actuación específico.

\begin{table}[H]
\centering
\caption{Tabla de riesgos y plan de contingencias}
\label{tab:riesgos}
\renewcommand{\arraystretch}{1.3}
\begin{tabularx}{\textwidth}{|X|X|}
\hline
\textbf{Riesgo identificado} & \textbf{Plan de contingencia} \\
\hline
Dificultad en la organización del tiempo durante el curso académico &
Establecer una planificación semanal realista y priorizar las tareas esenciales. \\
\hline
Retrasos en la implementación de algunos minijuegos &
Reducir la complejidad de ciertos minijuegos manteniendo el objetivo cognitivo principal. \\
\hline
Problemas técnicos durante el desarrollo en Unity &
Consultar la documentación oficial, recursos en línea y realizar pruebas frecuentes. \\
\hline
Sobrecarga de trabajo en la fase final del proyecto &
Adelantar la redacción de la memoria y reservar el último mes para revisión. \\
\hline
No completar los minijuegos previstos &
Priorizar al menos un minijuego representativo por categoría cognitiva. \\
\hline
Falta de validación con usuarios reales &
Realizar pruebas internas y justificar la validación limitada en el contexto académico. \\
\hline
Desajuste con los requisitos académicos &
Revisar periódicamente la guía del TFG y mantener comunicación con el tutor. \\
\hline
\end{tabularx}
\end{table}

\section{Análisis inicial de costes}

La valoración de la viabilidad económica de \textit{AttentiON} requiere
un desglose claro de los recursos utilizados y su impacto financiero real.
Al tratarse de un proyecto individual y académico, el enfoque principal ha
consistido en optimizar el uso de herramientas gratuitas o de bajo coste.

\begin{table}[H]
\centering
\caption{Resumen de costes del proyecto AttentiON}
\label{tab:costes}
\renewcommand{\arraystretch}{1.3}
\begin{tabularx}{\textwidth}{|l|X|r|}
\hline
\textbf{Categoría} & \textbf{Descripción} & \textbf{Coste (€)} \\
\hline
Software & Unity Personal (licencia gratuita para proyectos académicos) & 0{,}00 \\
\hline
Software & Visual Studio Code (editor de código, licencia gratuita) & 0{,}00 \\
\hline
Software & Assets gráficos con licencia libre o de uso académico & 0{,}00 \\
\hline
Hardware & Equipo informático propio del autor & 0{,}00 \\
\hline
Formación & Documentación oficial de Unity y recursos en línea & 0{,}00 \\
\hline
Coste de oportunidad & Horas de investigación, diseño y desarrollo
  (estimación: 400~h $\times$ 12~€/h) & 4.800{,}00 \\
\hline
\multicolumn{2}{|r|}{\textbf{Total directo}} & \textbf{0{,}00} \\
\hline
\multicolumn{2}{|r|}{\textbf{Total incluyendo coste de oportunidad}} &
  \textbf{4.800{,}00} \\
\hline
\end{tabularx}
\end{table}

El elemento que define la inversión real del proyecto es el \textbf{coste de
oportunidad}, dado por el tiempo personal dedicado a la investigación, la
programación y la redacción. Aunque no implica un movimiento de capital, esta
dedicación es la que ha transformado la propuesta teórica en un sistema
funcional.

\section{Impacto ambiental y responsabilidad social}

\subsection{Impacto ambiental}

El desarrollo de \textit{AttentiON} ha tenido un impacto ambiental reducido.
Al tratarse de un proyecto de software sin necesidad de producción de hardware
adicional, el consumo de recursos físicos se limita al uso del equipo informático
del autor durante el período de desarrollo.

El consumo energético asociado al uso de un ordenador de sobremesa o portátil
estándar durante el periodo de desarrollo (aproximadamente diez meses) se
estima en un rango razonable para un proyecto académico individual. No se han
generado residuos electrónicos, desplazamientos necesarios ni consumo de
materiales físicos relevante.

A futuro, la distribución digital de la plataforma como aplicación web o
de escritorio presentaría una huella ambiental significativamente inferior
a la de soluciones educativas basadas en materiales físicos (libros de texto,
materiales impresos, etc.).

\subsection{Responsabilidad social}

Desde la perspectiva de la responsabilidad social, \textit{AttentiON} tiene
un impacto potencialmente positivo. La plataforma está concebida para mejorar
las habilidades cognitivas de niños y niñas en edad escolar, lo que puede
contribuir a su desarrollo personal, a su rendimiento académico y a su
bienestar emocional.

En particular, se identifican las siguientes dimensiones de impacto social:

\begin{itemize}
  \item \textbf{Inclusión educativa:} al tratarse de una herramienta digital
    de acceso gratuito, puede facilitar el acceso a recursos de estimulación
    cognitiva a familias y centros educativos con presupuestos limitados.
  \item \textbf{Apoyo a colectivos con necesidades específicas:} la plataforma
    puede ser especialmente útil como recurso complementario para niños con
    dificultades de atención o de control inhibitorio.
  \item \textbf{Empoderamiento del profesorado:} los educadores disponen de
    una herramienta que enriquece su metodología sin requerir formación
    técnica especializada.
\end{itemize}

% ============================================================
%  CAPÍTULO 4 — METODOLOGÍA
% ============================================================
\chapter{Metodología}

\section{Justificación de la metodología}

Para desarrollar \textit{AttentiON} se ha adoptado una \textbf{metodología
iterativa e incremental}, cuya elección responde a la naturaleza práctica del
proyecto: un videojuego educativo que combina investigación teórica, diseño
de interacción y programación técnica de forma simultánea.

Este enfoque se adapta especialmente bien a proyectos de desarrollo de
software en los que los requisitos pueden evolucionar a lo largo del proceso.
Al dividir el trabajo en ciclos cortos, se posibilita la revisión continua de
lo desarrollado, la corrección temprana de errores y la mejora progresiva del
prototipo.

La elección se apoya, además, en precedentes del sector de los
\textit{serious games}: los videojuegos educativos de calidad raramente se
desarrollan de forma lineal; por el contrario, se benefician de procesos de
prototipado rápido, prueba y refinamiento iterativo.

\section{Fases del proyecto}

El desarrollo de \textit{AttentiON} se ha estructurado en cuatro fases principales,
si bien el proceso no ha sido estrictamente secuencial: la redacción de la memoria
y el desarrollo técnico han avanzado en paralelo durante la mayor parte del
proyecto.

\subsection{Fase 1: Diseño}

En esta fase se trabajó en la definición de las mecánicas, la estructura del
sistema y el catálogo de minijuegos. El objetivo principal fue diseñar una
interfaz intuitiva y accesible para el público infantil, garantizando la
alineación entre el componente visual y el funcional.

Los contenidos se organizaron en función de las capacidades cognitivas que se
buscaba estimular, estableciendo las cinco categorías que vertebran la
plataforma: atención, memoria, control de impulsos, gestión emocional y
planificación. Se determinó que cada categoría contaría con un mínimo de cinco
minijuegos distribuidos en tres niveles de dificultad.

\subsection{Fase 2: Desarrollo}

El desarrollo técnico se centró en la creación del prototipo dentro del motor
Unity. Se priorizó la implementación de los sistemas centrales (navegación,
gestión de escenas, sistema de dificultad, UI global) antes de abordar los
minijuegos individuales, dado que estos sistemas actúan como infraestructura
compartida.

Se adoptó un enfoque de desarrollo modular: cada minijuego se implementó como
una unidad independiente que hereda de la clase base \texttt{MinigameBase},
lo que facilitó la integración progresiva en el sistema global sin interrumpir
las partes ya funcionales.

\subsection{Fase 3: Pruebas y validación}

Una vez finalizadas las funcionalidades base, se llevaron a cabo pruebas
internas para verificar el funcionamiento del software y corregir errores de
programación. Estas pruebas sirvieron para ajustar las mecánicas y optimizar
la experiencia de usuario. Dado el marco académico del proyecto, la validación
se basó en revisiones técnicas propias y en pruebas de observación directa del
prototipo.

\subsection{Fase 4: Documentación y cierre}

De forma paralela a la programación, se elaboró la memoria del TFG, donde se
registraron todas las decisiones de diseño y los pasos seguidos. En esta fase
final se realizó la revisión completa del documento y se prepararon los archivos
necesarios para la entrega definitiva.

\section{Herramientas y tecnologías empleadas}

\begin{itemize}
  \item \textbf{Unity (versión 2022 LTS):} motor central del desarrollo.
    Elegido por su versatilidad, su extensa comunidad y por ser el estándar
    de la industria para videojuegos educativos. Permite centralizar en un
    mismo entorno el diseño de escenas, la gestión de elementos visuales y
    la programación, lo que ha sido clave para el desarrollo modular.
  \item \textbf{C\#:} lenguaje de programación utilizado para todos los
    scripts del proyecto. Su tipado estático y su integración nativa con
    Unity facilitan la detección temprana de errores.
  \item \textbf{Visual Studio Code:} editor de código ligero, con herramientas
    de depuración y soporte para C\# que agilizaron la corrección de errores
    durante las etapas de programación.
  \item \textbf{TextMesh Pro:} librería de Unity para el renderizado de texto
    de alta calidad, utilizada en todas las interfaces del juego.
  \item \textbf{Assets gráficos:} recursos procedentes de plataformas de
    assets digitales con licencias compatibles con el uso académico. Su
    utilización ha permitido optimizar el tiempo de producción.
  \item \textbf{Google Sheets:} herramienta de seguimiento de la planificación
    y elaboración del diagrama de Gantt.
  \item \textbf{Git / GitHub:} sistema de control de versiones del proyecto,
    que permite mantener el historial de cambios y recuperar versiones
    anteriores en caso de error.
\end{itemize}

\section{Seguimiento metodológico}

El sistema de trabajo iterativo se ha mantenido sin cambios significativos
durante todo el desarrollo. La decisión de abordar cada minijuego como una
unidad independiente antes de su integración definitiva ha demostrado ser
especialmente acertada: ha permitido detectar fallos en etapas tempranas
y ajustar las mecánicas de juego sin comprometer el sistema global.

Aunque durante el proceso fue necesario realizar pequeños ajustes, como la
reorganización de los minijuegos en categorías cognitivas o la concreción del
número de minijuegos por sección (fijado en cinco por categoría), estas
modificaciones se integraron de forma natural dentro del sistema de trabajo sin
necesidad de alterar la estrategia general.

El uso de ciclos cortos ha sido clave para gestionar el tiempo disponible,
priorizando la coherencia y la calidad del prototipo frente a un desarrollo
más extenso pero menos cuidado.

% ============================================================
%  CAPÍTULO 5 — DESARROLLO DEL PROYECTO
% ============================================================
\chapter{Desarrollo del proyecto}

Este capítulo describe el proceso técnico de desarrollo de
\textit{AttentiON}, desde la arquitectura global del sistema hasta la
implementación de cada categoría de minijuegos. Se justifican las decisiones
técnicas adoptadas, se describe el funcionamiento de los scripts y sistemas
principales, y se incluyen fragmentos de código representativos que ilustran
las soluciones implementadas.

\section{Arquitectura general del sistema}

El proyecto sigue una arquitectura en capas con responsabilidades bien
delimitadas:

\begin{enumerate}
  \item \textbf{Capa de núcleo (\textit{Core}):} gestores globales que
    persisten entre escenas y coordinan el estado del juego.
  \item \textbf{Capa de sistemas:} subsistemas especializados (UI, navegación,
    puntuación).
  \item \textbf{Capa de minijuegos:} implementaciones concretas de cada
    minijuego, con una clase base común.
  \item \textbf{Capa decorativa:} componentes visuales no funcionales.
\end{enumerate}

La estructura de carpetas del proyecto refleja directamente esta arquitectura:

\begin{verbatim}
Assets/_Project/Scripts/
  Core/
    GameManager.cs
    SceneLoader.cs
    DifficultyManager.cs
  Systems/
    UI/         (UIManager, IntroPanel, ButtonHoverScale, ...)
    Score/      (CoinManager)
    Navigation/ (LevelButton)
  Minigames/
    MinigameBase.cs
    Attention/
    Memory/
    ImpulseControl/
    EmotionalManagement/
    Planning/
  Decorative/   (CloudMover, Float, RotatingCoin)
\end{verbatim}

\begin{figure}[H]
  \centering
  \fbox{\parbox{12cm}{\centering\vspace{2cm}
    [Figura: Diagrama de arquitectura del sistema AttentiON]\vspace{2cm}}}
  \caption{Diagrama de arquitectura general del proyecto AttentiON.}
  \label{fig:arquitectura}
\end{figure}

\section{Sistemas del núcleo}

\subsection{GameManager — Singleton central}

\texttt{GameManager} es el gestor central de la aplicación. Implementa el
patrón \textit{Singleton} con persistencia entre escenas mediante
\texttt{DontDestroyOnLoad}, lo que garantiza que el estado global (dificultad
activa, puntuación acumulada) se conserve durante toda la sesión de juego.

Las responsabilidades de \texttt{GameManager} son:

\begin{itemize}
  \item Almacenar y exponer la dificultad seleccionada por el jugador.
  \item Mantener y actualizar la puntuación acumulada a través de la sesión.
  \item Proveer métodos de navegación global (al menú principal, al selector
    de dificultad).
\end{itemize}

\begin{lstlisting}[caption={GameManager.cs — Implementación del Singleton con persistencia},
                   label={lst:gamemanager}]
public class GameManager : MonoBehaviour
{
    public static GameManager Instance { get; private set; }

    private DifficultyLevel _currentDifficulty = DifficultyLevel.Easy;
    private int _totalScore = 0;

    public DifficultyLevel CurrentDifficulty => _currentDifficulty;
    public int TotalScore => _totalScore;

    [SerializeField] private string mainMenuScene   = "PrimeraPantalla";
    [SerializeField] private string difficultyScene = "DifficultySelector";

    private void Awake()
    {
        if (Instance != null && Instance != this)
        {
            Destroy(gameObject);
            return;
        }
        Instance = this;
        DontDestroyOnLoad(gameObject);
    }

    public void SetDifficulty(DifficultyLevel difficulty)
    {
        _currentDifficulty = difficulty;
    }

    public void AddScore(int amount)
    {
        _totalScore += amount;
    }

    public void ResetScore() => _totalScore = 0;

    public void GoToMainMenu()      => SceneLoader.LoadScene(mainMenuScene);
    public void GoToDifficultySelector() => SceneLoader.LoadScene(difficultyScene);
}

public enum DifficultyLevel
{
    Easy   = 0,
    Medium = 1,
    Hard   = 2
}
\end{lstlisting}

Los niveles de dificultad se modelan mediante el enumerado
\texttt{DifficultyLevel}, con tres valores que se corresponden con los
entornos temáticos de la plataforma: Bosque Mágico (6--8 años), Castillo
Volador (9--11 años) y Cueva Misteriosa (12+ años).

\subsection{SceneLoader — Sistema centralizado de navegación}

\texttt{SceneLoader} es una clase estática que centraliza toda la navegación
entre escenas del proyecto. Esta decisión de diseño aporta varias ventajas:
elimina las referencias a nombres de escena dispersas por el código, facilita
el refactoring y permite añadir lógica de transición (como efectos de fade)
en un único punto.

La clase define constantes para todos los nombres de escena del proyecto,
lo que evita errores tipográficos y hace el código más legible.

\begin{lstlisting}[caption={SceneLoader.cs — Constantes de escena y navegación adaptativa},
                   label={lst:sceneloader}]
public static class SceneLoader
{
    public const string MAIN_MENU           = "PrimeraPantalla";
    public const string DIFFICULTY_SELECTOR = "DifficultySelector";
    public const string GAME_SELECTOR_EASY  = "GameSelector";
    public const string GAME_SELECTOR_MEDIUM= "GameSelector2";
    public const string GAME_SELECTOR_HARD  = "GameSelector3";

    public static void LoadScene(string sceneName)
    {
        if (string.IsNullOrEmpty(sceneName)) return;
        SceneManager.LoadScene(sceneName);
    }

    public static void LoadGameSelector()
    {
        if (GameManager.Instance == null)
        {
            LoadScene(GAME_SELECTOR_EASY);
            return;
        }
        switch (GameManager.Instance.CurrentDifficulty)
        {
            case DifficultyLevel.Easy:   LoadScene(GAME_SELECTOR_EASY);   break;
            case DifficultyLevel.Medium: LoadScene(GAME_SELECTOR_MEDIUM); break;
            case DifficultyLevel.Hard:   LoadScene(GAME_SELECTOR_HARD);   break;
        }
    }

    public static void ReloadCurrentScene() =>
        LoadScene(SceneManager.GetActiveScene().name);

    public static void GoToMainMenu()
    {
        if (GameManager.Instance != null)
            GameManager.Instance.GoToMainMenu();
        else
            LoadScene(MAIN_MENU);
    }
}
\end{lstlisting}

El método \texttt{LoadGameSelector()} es especialmente relevante: consulta
la dificultad activa en \texttt{GameManager} y redirige al selector de juego
correspondiente, lo que implementa de forma transparente el sistema de tres
niveles.

\subsection{DifficultyManager — Gestor de selección de dificultad}

\texttt{DifficultyManager} es el componente que gestiona la pantalla de
selección de dificultad. Expone tres métodos públicos
(\texttt{SelectEasy()}, \texttt{SelectMedium()}, \texttt{SelectHard()}) que
se asignan directamente a los botones de la UI desde el Inspector de Unity.

Cuando el jugador selecciona una dificultad, el componente actualiza el
estado global en \texttt{GameManager} y navega automáticamente al selector
de juego apropiado.

\subsection{UIManager — Gestor global de interfaz}

\texttt{UIManager} complementa a \texttt{GameManager} con las
responsabilidades de interfaz. También implementa el patrón \textit{Singleton}
con persistencia entre escenas y gestiona:

\begin{itemize}
  \item El efecto de \textit{fade} (transición con fundido a negro) entre
    escenas.
  \item La visualización de la puntuación global.
  \item Los mensajes de estado temporales en pantalla.
\end{itemize}

El efecto de \textit{fade} se implementa mediante una corrutina que interpola
el canal alfa de un panel de imagen a pantalla completa, lo que proporciona
una transición suave sin necesidad de plugins externos.

\section{La clase base MinigameBase}

Uno de los pilares de la arquitectura del proyecto es la clase abstracta
\texttt{MinigameBase}, de la que heredan \textbf{todos} los minijuegos de la
plataforma. Esta decisión de diseño garantiza:

\begin{itemize}
  \item Comportamiento uniforme: todos los minijuegos muestran automáticamente
    un panel de introducción al cargarse.
  \item Interfaz contractual: los métodos abstractos
    (\texttt{OnMinigameStart}, \texttt{OnMinigameComplete},
    \texttt{OnMinigameFailed}) obligan a cada minijuego a implementar su
    lógica específica.
  \item Reutilización de código: métodos como \texttt{CompleteMinigame()},
    \texttt{FailMinigame()}, \texttt{ReturnToGameSelector()} o
    \texttt{RestartMinigame()} están disponibles para todas las subclases.
\end{itemize}

\begin{lstlisting}[caption={MinigameBase.cs — Clase base abstracta para todos los minijuegos},
                   label={lst:minigamebase}]
public abstract class MinigameBase : MonoBehaviour
{
    [SerializeField] protected string minigameName = "Minijuego";
    [SerializeField] protected MinigameCategory category = MinigameCategory.Memory;

    public int  Score     { get; protected set; } = 0;
    public bool IsPlaying { get; protected set; } = false;

    bool       _gameStarted;
    GameObject _introCvGO;

    protected virtual void Start()
    {
        BuildIntroPanel();
        StartCoroutine(WaitForSpaceKey());
    }

    IEnumerator WaitForSpaceKey()
    {
        while (!_gameStarted)
        {
            if (Input.GetKeyDown(KeyCode.Space)) LaunchGame();
            yield return null;
        }
    }

    void LaunchGame()
    {
        if (_gameStarted) return;
        _gameStarted = true;
        if (_introCvGO != null) Destroy(_introCvGO);
        IsPlaying = true;
        OnMinigameStart();
    }

    protected abstract void OnMinigameStart();
    protected abstract void OnMinigameComplete();
    protected abstract void OnMinigameFailed();

    protected void CompleteMinigame(int finalScore = 0)
    {
        if (!IsPlaying) return;
        IsPlaying = false;
        Score     = finalScore;
        if (GameManager.Instance != null)
            GameManager.Instance.AddScore(finalScore);
        OnMinigameComplete();
    }

    protected void FailMinigame()
    {
        if (!IsPlaying) return;
        IsPlaying = false;
        OnMinigameFailed();
    }

    protected void ReturnToGameSelector() => SceneLoader.LoadGameSelector();
    protected void RestartMinigame()      => SceneLoader.ReloadCurrentScene();
}

public enum MinigameCategory
{
    Memory = 0, ImpulseControl = 1, EmotionalManagement = 2,
    Attention = 3, Planning = 4
}
\end{lstlisting}

El panel de introducción (\texttt{IntroPanel}) se construye de forma
procedural en tiempo de ejecución: se genera un canvas completo con título,
categoría, descripción de la mecánica y botón de inicio. El código de color
del panel se asigna automáticamente según la categoría del minijuego, lo
que refuerza la coherencia visual del sistema.

\section{Categoría: Atención}

Los minijuegos de atención están diseñados para entrenar la capacidad de
concentración, la atención selectiva y la velocidad de respuesta ante
estímulos. Se han implementado tres minijuegos principales en esta categoría.

\subsection{Reacción Rápida (QuickReactionGameManager)}

Este minijuego mide el tiempo de reacción del jugador ante estímulos visuales.
La mecánica consiste en una serie de rondas en las que el jugador debe hacer
clic en el momento exacto en que aparece un estímulo verde, evitando pulsar
antes de tiempo o dejar que el tiempo expire.

El sistema gestiona tres tipos de resultado por ronda:
\begin{enumerate}
  \item \textbf{Acierto}: el jugador hace clic dentro del límite de tiempo.
  \item \textbf{Anticipación}: el jugador pulsa antes de que aparezca el
    estímulo.
  \item \textbf{Tiempo agotado}: el jugador no responde a tiempo.
\end{enumerate}

Los parámetros de dificultad (tiempo de espera mínimo y máximo antes del
estímulo, número de rondas, límites de tiempo de reacción por ronda) son
completamente configurables desde el Inspector de Unity, lo que facilita
el ajuste para cada nivel de dificultad.

\subsection{Seguimiento de Objeto (TrackingGameManager)}

En este minijuego el jugador debe seguir visualmente un objeto en movimiento
durante un período determinado. El componente \texttt{TrackingDetector}
detecta si el cursor del jugador se mantiene dentro del área del objeto
durante el tiempo requerido.

La mecánica trabaja específicamente la atención sostenida: el jugador no debe
reaccionar a un estímulo puntual, sino mantener el foco en un elemento en
continuo movimiento, simulando tareas de vigilancia visual.

\subsection{Cambio de Regla (RuleSwitchGameManager)}

Este minijuego entrena la flexibilidad cognitiva mediante la presentación de
estímulos que deben clasificarse según una regla que cambia periódicamente.
El sistema gestiona tres subsistemas coordinados: el
\texttt{RuleSwitchRuleManager} (gestión de reglas), el
\texttt{RuleSwitchStimulusManager} (generación de estímulos) y el
\texttt{RuleSwitchUIController} (interfaz).

\begin{figure}[H]
  \centering
  \fbox{\parbox{12cm}{\centering\vspace{2cm}
    [Figura: Captura de pantalla de los minijuegos de Atención]\vspace{2cm}}}
  \caption{Vista de los minijuegos de la categoría Atención en AttentiON.}
  \label{fig:atencion}
\end{figure}

\section{Categoría: Memoria}

Los minijuegos de memoria están orientados al entrenamiento de la memoria
operativa y la memoria visual. Se han implementado varios minijuegos con
mecánicas diversas para mantener el interés del jugador.

\subsection{Repite el Dibujo (PatternGameController)}

Este minijuego presenta al jugador un patrón de celdas activadas en una
cuadrícula, que debe memorizar durante un tiempo limitado y reproducir
correctamente a continuación. Consta de tres rondas progresivas que aumentan
la complejidad de la cuadrícula y reducen el tiempo de memorización.

\begin{lstlisting}[caption={PatternGameController.cs — Configuración de rondas},
                   label={lst:pattern}]
public class PatternGameController : MinigameBase
{
    [Header("Ronda 1")]
    public int   colsRound1    = 3;
    public int   rowsRound1    = 3;
    public int   patternRound1 = 4;
    public float timeRound1    = 4f;

    [Header("Ronda 2")]
    public int   colsRound2    = 4;
    public int   rowsRound2    = 4;
    public int   patternRound2 = 6;
    public float timeRound2    = 3f;

    [Header("Ronda 3")]
    public int   colsRound3    = 5;
    public int   rowsRound3    = 5;
    public int   patternRound3 = 9;
    public float timeRound3    = 2f;

    private int  _currentRound = 0;
    private enum Phase { Memorize, Recall, Result }
    private Phase _phase;
    // ...
}
\end{lstlisting}

El ciclo de juego de cada ronda se divide en dos fases claramente diferenciadas:
\textbf{Memorizar} (el patrón se muestra durante el tiempo configurado) y
\textbf{Reproducir} (el jugador selecciona las celdas que recuerda). Si el
jugador falla cualquier ronda, el juego termina con derrota; si completa las
tres rondas, se registra la victoria.

\subsection{Cambios Sutiles (FindChangeGameManager)}

Este minijuego muestra al jugador una escena con formas de colores, la retira
brevemente mediante un fundido a negro y la reemplaza por una versión
ligeramente modificada. El jugador debe identificar el elemento que ha cambiado.

El flujo de juego se gestiona mediante un enumerado de fases
(\texttt{Observe}, \texttt{Transition}, \texttt{Find}, \texttt{Result})
y corrutinas encadenadas, lo que permite controlar con precisión los tiempos
de cada etapa.

\subsection{Simon Says}

Adaptación digital del clásico juego de memoria secuencial. El sistema muestra
una secuencia de colores que se amplía en cada ronda. El jugador debe reproducir
la secuencia exacta en el mismo orden. El componente \texttt{SimonGame} gestiona
la generación de secuencias, la reproducción de la animación de muestra y la
validación de la entrada del jugador.

\subsection{Memoria de Posiciones (PositionMemoryGameManager)}

El jugador debe memorizar la posición de una serie de objetos que se muestran
brevemente y, a continuación, identificar cuál de ellos ha cambiado de lugar.
La dificultad aumenta con el número de objetos y la reducción del tiempo de
observación.

\subsection{Palabras Fugaces (WordMemoryGameManager)}

Minijuego de memoria verbal: se muestran palabras durante un tiempo limitado
y el jugador debe identificarlas posteriormente entre un conjunto de
distractores. Entrena la memoria semántica y la atención selectiva a estímulos
verbales.

\begin{figure}[H]
  \centering
  \fbox{\parbox{12cm}{\centering\vspace{2cm}
    [Figura: Captura de pantalla de los minijuegos de Memoria]\vspace{2cm}}}
  \caption{Vista de los minijuegos de la categoría Memoria en AttentiON.}
  \label{fig:memoria}
\end{figure}

\section{Categoría: Control de impulsos}

Los minijuegos de control de impulsos buscan entrenar la capacidad de inhibir
respuestas automáticas y actuar de forma reflexiva. Esta es una de las
habilidades ejecutivas más críticas en el desarrollo infantil.

\subsection{No Pulses Todavía (DontPressGameManager)}

Este minijuego pone a prueba directamente la capacidad de inhibición del
jugador: se presenta un botón de gran tamaño que el jugador debe
\textbf{no pulsar} durante un período de tiempo que va aumentando
progresivamente. Cualquier pulsación prematura supone el fallo de la ronda.

El \texttt{DontPressTimerManager} gestiona el temporizador y la presión
creciente sobre el jugador, mientras que el \texttt{DontPressUIController}
proporciona señales visuales que aumentan la tensión de forma deliberada.

\subsection{Respuesta Inversa (InverseResponseGameManager)}

En este minijuego el jugador recibe estímulos visuales y debe responder de
forma \textbf{contraria} a la respuesta intuitiva. Por ejemplo, si aparece
una flecha hacia la derecha, el jugador debe pulsar el botón de la izquierda.
Esta mecánica, inspirada en la tarea Go/No-Go de la psicología cognitiva,
entrena de forma efectiva el control inhibitorio.

\begin{lstlisting}[caption={InverseResponseGameManager.cs — Mecánica de respuesta inversa},
                   label={lst:inverse}]
public class InverseResponseGameManager : MinigameBase
{
    protected override void OnMinigameStart()
    {
        _input    = GetComponent<InverseResponseInputHandler>();
        _stimulus = GetComponent<InverseResponseStimulusManager>();
        _ui       = GetComponent<InverseResponseUIController>();

        _stimulus.OnStimulusGenerated += HandleStimulus;
        _input.OnResponseReceived     += HandleResponse;

        StartNextRound();
    }

    void HandleResponse(bool pressedLeft)
    {
        bool stimulusIsLeft = _currentStimulus.IsLeft;
        bool correct = pressedLeft != stimulusIsLeft;

        if (correct) RegisterCorrect();
        else         RegisterError();
    }
    // ...
}
\end{lstlisting}

\subsection{Stop \& Go (StopAndGoGameManager)}

Minijuego de reacción selectiva: objetos de distintos tipos se mueven por la
pantalla y el jugador debe detener únicamente los que cumplen una condición
específica (por ejemplo, solo los de color rojo), ignorando los demás. El
componente \texttt{StopAndGoZoneManager} gestiona las zonas de detección, y
\texttt{StopAndGoObjectMover} controla el movimiento de los objetos.

\begin{figure}[H]
  \centering
  \fbox{\parbox{12cm}{\centering\vspace{2cm}
    [Figura: Captura de pantalla de los minijuegos de Control de impulsos]\vspace{2cm}}}
  \caption{Vista de los minijuegos de la categoría Control de impulsos en AttentiON.}
  \label{fig:impulsos}
\end{figure}

\section{Categoría: Gestión emocional}

Los minijuegos de gestión emocional buscan que el jugador identifique,
comprenda y regule emociones en situaciones simuladas. Esta categoría conecta
con la dimensión afectiva del aprendizaje, frecuentemente poco trabajada en
plataformas de entrenamiento cognitivo.

\subsection{Consecuencias Emocionales (ConsequencesGameManager)}

Se presentan situaciones cotidianas al jugador, que debe elegir la reacción
emocional más adecuada entre varias opciones. El sistema registra si la
respuesta es positiva, neutra o negativa, y asigna puntuación en consecuencia.

\begin{lstlisting}[caption={ConsequencesGameManager.cs — Evaluación de respuestas emocionales},
                   label={lst:consequences}]
public class ConsequencesGameManager : MinigameBase
{
    public int situationCount = 5;
    public int roundsToWin    = 3;
    public int pointsPositive = 20;
    public int pointsNeutral  = 8;
    public int pointsNegative = 0;

    void HandleOptionChosen(EmotionalResponse response)
    {
        int pts = response switch
        {
            EmotionalResponse.Positive => pointsPositive,
            EmotionalResponse.Neutral  => pointsNeutral,
            _                          => pointsNegative
        };
        _score += pts;
        if (response == EmotionalResponse.Positive) _positiveCount++;
        _ui.ShowConsequence(response);
    }

    void EndGame()
    {
        bool won = _positiveCount >= roundsToWin;
        if (won) CompleteMinigame(_score);
        else     FailMinigame();
    }
}
\end{lstlisting}

\subsection{Regulación Progresiva (RegulationGameManager)}

Este minijuego introduce al jugador en una situación de nivel emocional
elevado que debe reducir aplicando técnicas de regulación (respiración,
caminar, hablar, etc.). La mecánica clave es que el nivel emocional aumenta
automáticamente cada turno, de forma que las acciones menos eficaces no
compensan la tensión acumulada. Esto obliga al jugador a tomar decisiones
estratégicas sobre qué técnica aplicar y cuándo.

\subsection{Equilibrio Emocional (EmotionalBalanceController)}

Minijuego de equilibrio dinámico: el jugador debe mantener un indicador
dentro de una zona segura respondiendo a eventos emocionales positivos y
negativos. La mecánica entrena la regulación emocional como proceso continuo
en lugar de como respuesta puntual.

\subsection{Control de Atracción (AttractionGameManager)}

El jugador controla un cursor que debe evitar ser atraído por ciertos
elementos perturbadores en pantalla. La atracción simulada representa
distracciones emocionales, y el jugador debe ejercer un esfuerzo deliberado
para mantenerse en la trayectoria objetivo.

\begin{figure}[H]
  \centering
  \fbox{\parbox{12cm}{\centering\vspace{2cm}
    [Figura: Captura de pantalla de los minijuegos de Gestión emocional]\vspace{2cm}}}
  \caption{Vista de los minijuegos de la categoría Gestión emocional en AttentiON.}
  \label{fig:emocional}
\end{figure}

\section{Categoría: Planificación}

Los minijuegos de planificación exigen que el jugador anticipe consecuencias,
organice acciones en secuencia y gestione recursos de forma eficiente. Esta
categoría trabaja las habilidades ejecutivas de orden superior.

\subsection{Secuencia de Acciones (ActionSequenceController)}

El jugador debe ordenar correctamente una secuencia de acciones para completar
una tarea cotidiana (por ejemplo: levantarse, desayunar e ir al colegio). Los
botones se presentan en orden aleatorio y el jugador debe activarlos en la
secuencia correcta.

El minijuego consta de tres rondas progresivas con secuencias de diferente
complejidad. Un error en cualquier paso reinicia la ronda desde el principio,
lo que refuerza la importancia del orden y la planificación previa.

\subsection{Gestión de Recursos (ResourceGameController)}

Este minijuego propone al jugador un escenario de gestión de recursos: dispone
de una cantidad limitada de «estrellas» (energía) y debe alcanzar un objetivo
de progreso eligiendo acciones, cada una con un coste y un beneficio diferente.
La dificultad reside en optimizar el gasto para alcanzar el objetivo sin
quedarse sin recursos.

\begin{lstlisting}[caption={ResourceGameController.cs — Gestión de acciones con coste y progreso},
                   label={lst:resource}]
[Serializable]
public class ActionData
{
    public string icon       = "★";
    public string actionName = "Accion";
    public int    cost       = 1;
    public int    progress   = 10;
    public Color  buttonColor = Color.black;
}

public class ResourceGameController : MinigameBase
{
    public int starsEasy   = 20;
    public int starsMedium = 14;
    public int starsHard   = 9;
    public int goal        = 100;

    public List<ActionData> actions = new List<ActionData>();

    void ExecuteAction(ActionData action)
    {
        if (_stars < action.cost) return;
        _stars   -= action.cost;
        _progress += action.progress;
        RefreshUI();
        if (_progress >= goal) CompleteMinigame(CalculateScore());
    }
}
\end{lstlisting}

\subsection{Ruta Óptima (OptimalPathController)}

El jugador debe encontrar el camino más eficiente entre un punto de inicio y
un destino en un mapa con obstáculos y restricciones. El sistema evalúa no
solo si el jugador llega al destino, sino si lo hace con el menor número de
pasos posible.

\subsection{Orden Correcto (OrderGameController)}

Minijuego de secuenciación visual: se muestran elementos desordenados que el
jugador debe organizar siguiendo un criterio específico (numérico, temporal,
lógico). El componente \texttt{OrderManager} valida cada paso y proporciona
retroalimentación inmediata.

\begin{figure}[H]
  \centering
  \fbox{\parbox{12cm}{\centering\vspace{2cm}
    [Figura: Captura de pantalla de los minijuegos de Planificación]\vspace{2cm}}}
  \caption{Vista de los minijuegos de la categoría Planificación en AttentiON.}
  \label{fig:planificacion}
\end{figure}

\section{Sistema de dificultad}

El sistema de dificultad adaptativa de \textit{AttentiON} opera en dos niveles:

\begin{enumerate}
  \item \textbf{Nivel de sesión:} el jugador selecciona una dificultad global
    al inicio de cada sesión, que determina a qué \textit{GameSelector} se
    dirige y qué escenas de minijuegos se cargan (sufijos \texttt{\_Easy},
    \texttt{\_Medium}, \texttt{\_Hard} en los nombres de escena).
  \item \textbf{Nivel interno de minijuego:} cada minijuego expone sus
    parámetros de dificultad en el Inspector de Unity, lo que permite ajustar
    el tiempo de reacción, el número de rondas, la velocidad de los objetos u
    otros parámetros específicos para cada escena.
\end{enumerate}

Los tres niveles temáticos (\textit{Bosque Mágico}, \textit{Castillo Volador},
\textit{Cueva Misteriosa}) no son únicamente nombres: cada uno va acompañado
de una estética visual diferenciada y se corresponde con una franja de edad
objetivo, lo que refuerza la identidad de la plataforma y orienta tanto a los
jugadores como a los educadores sobre el uso apropiado.

\section{Sistema de UI e IntroPanel}

La clase estática \texttt{IntroPanel} genera proceduralmente, mediante código C\#,
un canvas completo de introducción para cualquier minijuego que herede de
\texttt{MinigameBase}. Esta solución evita la dependencia de prefabs externos y
garantiza una presentación coherente en toda la plataforma.

Cada panel muestra el nombre del minijuego, la categoría a la que pertenece
(con el color asociado), la descripción de la mecánica y un botón de inicio.
El sistema también detecta si hay un \texttt{EventSystem} activo y lo crea si
no existe, garantizando que la interacción con la UI funcione correctamente en
cualquier escena.

\begin{table}[H]
\centering
\caption{Colores de categoría en el sistema de IntroPanel}
\label{tab:colores}
\begin{tabular}{ll}
\toprule
\textbf{Categoría} & \textbf{Color} \\
\midrule
Memoria              & Púrpura \texttt{(0.58, 0.28, 0.92)} \\
Control de impulsos  & Naranja \texttt{(0.95, 0.55, 0.12)} \\
Gestión emocional    & Verde   \texttt{(0.18, 0.80, 0.58)} \\
Atención             & Amarillo \texttt{(0.98, 0.80, 0.10)} \\
Planificación        & Azul    \texttt{(0.28, 0.60, 1.00)} \\
\bottomrule
\end{tabular}
\end{table}

\section{Elementos decorativos}

Los scripts de la carpeta \texttt{Decorative} aportan vida visual a los menús
y pantallas de la plataforma sin influir en la lógica del juego:

\begin{itemize}
  \item \texttt{CloudMover}: desplaza objetos con forma de nube a lo largo
    de la pantalla y los reposiciona al salir del área visible, creando un
    efecto de cielo en movimiento en los menús.
  \item \texttt{Float}: aplica un movimiento senoidal vertical a un objeto,
    generando la sensación de flotación.
  \item \texttt{RotatingCoin}: hace rotar continuamente un objeto 3D,
    utilizado como indicador visual del sistema de monedas.
\end{itemize}

\section{Sistema de puntuación}

El sistema de puntuación opera en dos capas:

\begin{itemize}
  \item \textbf{Puntuación de sesión:} gestionada por \texttt{GameManager},
    que acumula los puntos obtenidos en todos los minijuegos de la sesión
    actual mediante el método \texttt{AddScore()}.
  \item \textbf{Sistema de monedas:} \texttt{CoinManager} implementa un
    sistema de recompensas por minijuego que actúa como retroalimentación
    inmediata para el jugador. El número de monedas recogidas se muestra
    en pantalla y motiva al jugador a completar cada minijuego con el
    mayor número posible de aciertos.
\end{itemize}

\section{Resumen de minijuegos implementados}

\begin{table}[H]
\centering
\caption{Resumen de minijuegos implementados en AttentiON}
\label{tab:minijuegos}
\renewcommand{\arraystretch}{1.3}
\begin{tabularx}{\textwidth}{|l|X|l|}
\hline
\textbf{Categoría} & \textbf{Minijuego} & \textbf{Habilidad principal} \\
\hline
\multirow{5}{*}{Atención}
  & Reacción Rápida          & Velocidad de respuesta \\
  & Seguimiento de Objeto    & Atención sostenida \\
  & Cambio de Regla          & Flexibilidad cognitiva \\
  & (2 minijuegos adicionales en desarrollo) & \\
\hline
\multirow{5}{*}{Memoria}
  & Repite el Dibujo          & Memoria visual \\
  & Cambios Sutiles           & Memoria de trabajo \\
  & Simon Says                & Memoria secuencial \\
  & Memoria de Posiciones     & Memoria espacial \\
  & Palabras Fugaces          & Memoria verbal \\
\hline
\multirow{4}{*}{Control de impulsos}
  & No Pulses Todavía         & Inhibición de respuesta \\
  & Respuesta Inversa         & Control inhibitorio \\
  & Stop \& Go                & Atención selectiva \\
  & (2 minijuegos adicionales en desarrollo) & \\
\hline
\multirow{4}{*}{Gestión emocional}
  & Consecuencias Emocionales & Toma de decisiones \\
  & Regulación Progresiva     & Autorregulación \\
  & Equilibrio Emocional      & Regulación continua \\
  & Control de Atracción      & Gestión de distracciones \\
\hline
\multirow{4}{*}{Planificación}
  & Secuencia de Acciones     & Planificación temporal \\
  & Gestión de Recursos       & Razonamiento estratégico \\
  & Ruta Óptima               & Anticipación y toma de decisiones \\
  & Orden Correcto            & Secuenciación lógica \\
\hline
\end{tabularx}
\end{table}

\section{Decisiones técnicas clave}

A lo largo del desarrollo se han tomado varias decisiones de diseño técnico
relevantes que conviene justificar:

\begin{itemize}
  \item \textbf{Uso del patrón Singleton con DontDestroyOnLoad:} esta elección
    garantiza que el estado global (dificultad, puntuación) persiste entre
    escenas sin necesidad de sistemas de serialización complejos, lo que resulta
    adecuado para la escala del proyecto.
  \item \textbf{Generación procedural de UI:} la construcción por código de
    paneles como \texttt{IntroPanel} elimina la dependencia de prefabs y facilita
    la evolución del sistema sin riesgo de inconsistencias visuales entre escenas.
  \item \textbf{Clase base abstracta MinigameBase:} la herencia obliga a cada
    minijuego a implementar los métodos del contrato, reduciendo el riesgo de
    olvidar implementar comportamientos críticos (como registrar el resultado en
    \texttt{GameManager}).
  \item \textbf{Separación de responsabilidades dentro de cada minijuego:}
    los minijuegos más complejos dividen su lógica en componentes hermanos
    (GameManager, UIController, InputHandler, StimulusManager, etc.),
    siguiendo el principio de responsabilidad única y facilitando el
    mantenimiento y la extensión.
  \item \textbf{Centralización de nombres de escena en SceneLoader:}
    evita la dispersión de strings literales por el código, que es una fuente
    habitual de errores difíciles de detectar en proyectos con muchas escenas.
\end{itemize}

% ============================================================
%  BIBLIOGRAFÍA
% ============================================================
{\setlength{\bibsep}{4pt}
\bibliographystyle{apalike}
\bibliography{referencias}
\addcontentsline{toc}{chapter}{Bibliografía}}

\end{document}
